problem:  
Let \( (X^8, g) \) be a compact Riemannian manifold with torsion-free \( \mathrm{Spin}(7)\)-structure characterized by a parallel 4-form \( \Phi \).  

Assume \(X\) admits a cohomogeneity-one isometric action by a compact, connected Lie group \(G\) that preserves \( \Phi \), so that the generic orbits are 7-dimensional homogeneous spaces \(G/H\).  In a tubular neighborhood of principal orbits, the metric may be expressed as  
\[
g = dr^2 + g_r,
\]
where each \(g_r\) is a smooth \(G\)-invariant metric on \(G/H\).  The induced family of invariant \(G_2\)-structures \( \varphi_r \) on \(G/H\) satisfies the **Hitchin flow equations** derived from the condition that the holonomy of \(g\) is contained in \( \mathrm{Spin}(7)\).  

Known examples (Bryant–Salamon metrics) yield complete noncompact manifolds via such Hitchin flow evolution, starting from nearly-parallel \(G_2\)-structures on spheres or certain 3-Sasakian spaces.  

**Conjectural Problem:**  
Does there exist a compact torsion-free \(\mathrm{Spin}(7)\) manifold \( (X, \Phi) \) admitting a cohomogeneity-one \(G\)-action preserving \(\Phi\), other than the flat torus?  

Equivalently, can one find a \(G\)-invariant family of \(G_2\)-structures \(\varphi_r\) on \(G/H\) satisfying Hitchin’s flow equations on a **finite interval** \(r\in [r_-, r_+]\), such that the metric closes smoothly at each endpoint—possibly through the appearance of **singular orbits** of types \(G/H_-\) and \(G/H_+\)—yielding a smooth compact space \(X\)?  

Moreover, could such a solution be realized by gluing two asymptotically conical noncompact \(\mathrm{Spin}(7)\) manifolds along a common principal orbit \(G/H\), providing new topological types of compact \(\mathrm{Spin}(7)\) manifolds not encompassed by Joyce’s desingularization constructions?

Why is it a "good" problem:  
This problem targets a fundamental open question at the intersection of **exceptional holonomy geometry**, **group actions**, and **geometric analysis**: whether symmetry reduction via cohomogeneity‑one actions can ever produce compact manifolds of Spin(7) holonomy.  

It bridges multiple domains:  
- **Analytic:** the existence of smooth finite-interval solutions of Hitchin’s nonlinear flow with prescribed boundary singularities.  
- **Geometric/topological:** the classification of possible smooth closings and singular orbit types compatible with holonomy \( \mathrm{Spin}(7)\).  
- **Constructive:** the potential for gluing noncompact exceptional metrics into compact ones.  

The statement is simple—asking whether compact symmetric Spin(7) manifolds exist—but its resolution demands deep analytic insight into geometric flows, Lie group representation, and boundary-value geometry. A positive solution would produce entirely new compact examples; a negative one would reveal strong rigidity in the presence of symmetry.  

Thus, the problem is conceptually clean, mathematically sound, and unifies critical ideas across holonomy theory, differential geometry, and nonlinear analysis, meeting the standard of a “good” research problem.